\selectlanguage{english}\Simplified
\Language{zh-Hans}{简体中文}{使用指南}
\firsttopic{启动、休眠与返回笔记}
Inkline 在 reMarkable 2 上运行本地 shell。使用 Type Folio 或 USB 键盘，可以在平板上
工作，也可以通过 SSH 或 Goblin Mosh 连接其他计算机。文字和图像以适合电子纸的灰度显示。

\begin{notice}
按 \key{Opt+RightAlt+T}，即可直接在平板上打开 Inkline。\par
\key{电源键：}按下再松开即可休眠，再按一次唤醒。已打开的 shell 和编辑器缓冲区保留在 RAM 中。
\end{notice}
休眠会暂停设备，不会退出 Inkline，也不会把设置写入闪存。唤醒后，网络连接可能需要重新建立。
用于唤醒的按键不会再次触发休眠。更新后，请退出并重新打开 Inkline 一次，以启用新的电源键处理。

要返回笔记，点按 \key{Quit} 或按 \key{Ctrl+Shift+Q}，然后确认。
输入 \code{exit} 只关闭当前终端；关闭最后一个终端后返回笔记。
底栏包含 Escape、历史记录上／下、Quit 和 Settings 按钮。
启动时的简短说明和小 Goblin 标志可用 \code{clear} 清除。

如果快捷键启动的应用占用了屏幕，按 \key{Opt+RightAlt+T} 返回 Inkline。
紧急恢复时，\key{按住 Opt+RightAlt+Backspace 两秒}。
此操作会关闭所有 Inkline 终端，丢失这些会话中尚未保存的工作。
设备重启后仍默认进入笔记界面。

RightAlt 指空格键右侧的 Alt/Opt 键。按住它和独立的 Opt 键，再按 T；不要按 Ctrl。单独使用 Super 的组合仍可自定义。

\ProjectLinks

\topic{按键与终端槽位}
Type Folio 在 Ctrl 与 Alt 之间有一个独立的 \key{Opt} 键，另有一个
\key{右 Alt/Opt} 键。两者用途不同。

\begin{keytable}
\key{按键} & \key{操作} \\ \midrule
独立 Opt + 1--0 & F1--F10 \\
右 Alt/Opt + 1--9 & 选择终端槽位1--9 \\
右 Alt/Opt + Left / Right & 上一个／下一个槽位 \\
右 Alt/Opt + Tab & Escape \\
右 Alt/Opt + Up / Down & PageUp / PageDown \\
右 Alt/Opt + Backspace & 确认退出所有终端 \\
任一 Alt + Space & Unicode 键盘 \\
右 Alt/Opt + C / V & 复制选区／粘贴 \\
Ctrl+Shift+B & 隐藏／显示底栏 \\
\end{keytable}

在\key{美式 Type Folio} 上，右 Alt/Opt+\key{0} 输入 \code{+}；
右 Alt/Opt+\key{紧挨 Backspace 左侧的减号键}输入 \code{=}。
关闭输入法后，Shift+6 直接输入插入符号：\code{c^2} 保持为 \code{c^2}。
没有用于缩放字体的键盘组合键。

USB 键盘可直接使用普通功能键。独立 Opt 的快捷操作在提供 Meta 的 USB 键盘上使用 Meta 修饰键。

最多有\key{九个独立槽位}。选择空槽位会打开一个 shell。指示器显示当前打开的终端数，而非容量。
如果只有槽位1和9打开，在9中显示 \key{Terminal 2 / 2 · slot 9}。
关闭 shell 会释放槽位，不会重新编号其他快捷键。

每个终端都有自己的 shell、工作目录、输入组合状态、图像和500行回滚历史。
显示设置和输入法设置对所有槽位生效。

\topic{设置与显示}
点按底栏第五个按钮 \key{Settings}。如果底栏隐藏，按 \key{Alt+Space}，再按
\key{F2} 或点按 Unicode 键盘中的 Settings。实心选项表示当前状态，虚线边框表示键盘焦点。
Tab 移动焦点，方向键和 Enter 更改选择。

\key{Caps Lock：}可选择 Control（默认）或 Caps Lock。
\key{字体大小：}双指捏合，或使用设置中的减号／加号。
缓慢移动可每次调整一个像素，范围为\key{6至48像素}。

\key{Text darkness：}50\% 为原始渲染，数值越高，字形边缘越深。
\key{Minimum contrast} 在应用选择相近颜色时拉开前景与背景的亮度差，默认35\%。

\begin{choices}
\key{电子纸模式} & \key{取舍} \\ \midrule
Fast & 合并等待最少，使用动画波形，残影较多 \\
Balanced & UI 波形，按32 ms 合并输出 \\
Crisp（默认） & 灰度更清晰，使用内容波形，更新较慢 \\
Mono & 快速黑白显示，舍弃图像中的灰色层次 \\
Saver & 按120 ms 合并，减少连续输出的刷新次数 \\
\end{choices}

只重绘发生变化的文字行，但面板本身仍需要刷新时间。
Saver 会减少刷新请求，省电效果尚未测量。升级会保留明确保存的模式，不会强制覆盖为新默认值。

\begin{notice}
所有设置在使用期间保留在 \key{RAM}，正常退出 Inkline 时一起保存。
抬起手指、关闭设置和休眠都不会写入闪存。没有实际更改的会话不会写入设置。
崩溃、强制终止或断电可能丢失尚未保存的设置更改。
\end{notice}

\topic{Unicode 与输入法}
按 \key{Alt+Space} 打开 Unicode 键盘。点按字符即可输入，键盘保持打开，方便连续输入。
点按 \key{Done} 或按 Escape 关闭。

分类包括字母、组合标记、数字、标点、数学／货币、符号、空格、格式字符和私用区。
上下拖动滚动；也可以使用方向键、PageUp/PageDown、Home/End 或鼠标滚轮。
\key{Tab} 切换分类，Enter 输入选中的字符。

输入十六进制 Unicode 值（例如 \code{03bb}）后按 Enter，可以输入 λ。
Backspace 编辑数值。组合标记旁的虚线圆仅用于展示，实际只输入标记。
空格和格式字符带有可见标签。字符表使用 Qt 的 Unicode 数据，不包括控制字符、
未分配码位、代理码位和非字符。部分字符需要未安装的字体，私用区需要匹配的字体。

在键盘中按 \key{F2} 或点按 \key{Settings}，然后选择 \key{Input method}：
\begin{choices}
Off — direct keyboard & 普通美式键盘直接输入 \\
Romaji & 从罗马字输入日语平假名 \\
Pinyin & 拼音输入中文 \\
Zhuyin & 使用基本注音键位输入中文 \\
Wubi 86 & 使用五笔86字形编码输入中文 \\
US-International & 常用拉丁字母重音 \\
\end{choices}

在候选栏选择结果。要恢复普通美式输入，在同一菜单选择
\key{Off — direct keyboard}。此选择对所有终端生效，并在 Inkline 退出时保存。

\topic{触控、手写笔与剪贴板}
\key{双指上下滑动：}滚动历史。向下拖动查看较早输出，向上返回提示符方向。
每个终端除当前屏幕外，还在 RAM 中保留500个物理行。

\key{双指左右滑动：}每次手势切换一个终端槽位。再次切换前请抬起手指。
双指张开或合拢可调整字体大小。

\key{单指：}拖动选择文字。右 Alt/Opt+C 复制，右 Alt/Opt+V 粘贴。
所有终端共享一个最大128 KiB 的 RAM 剪贴板。

\key{手写笔：}启用终端鼠标报告的应用会收到对应的鼠标左键按下、移动与松开事件，
使用应用协商的协议，包括 SGR。其他情况下拖动用于选择文字。
按住 \key{Shift}，可以在启用鼠标报告的应用中强制进行本地选择。

\key{OSC 52} 允许终端程序（包括经 SSH 运行的远程程序）读取或替换 Inkline 的 RAM 剪贴板。
它与笔记应用的剪贴板不同。右 Alt/Opt+X 会复制，并提示在编辑器内部执行删除；
OSC 52 无法删除编辑器里的文字。

\key{图像：}kitty 的内联与共享内存传输保留在 RAM 中；文件与临时文件传输禁用。
显式请求时支持 sixel，但尚未自动声明支持。kitty 占位符、动画播放和部分 sixel 模式尚未完成。

\code{clear} 和全屏擦除会同时清除可见图像与文字。内存不足时会回收屏幕外的图像。
没有磁盘图像缓存，但大图仍会占用设备有限的 RAM。

\topic{USB 键盘与打字机}
打开\key{Settings → USB（第 3 页）}。\key{Network} 恢复 USB 以太网；
\key{Send keyboard} 让平板成为另一台电脑的键盘；\key{Receive keyboard}
通过有供电的 OTG 转接器接收有线键盘。Type Folio 在各模式下均可使用。
输出键盘模式需要先安装网站提供的 Python 工具，无需额外的 Python 模块。

选择 Send keyboard，用 USB-C 连接电脑，设置下列接收端配置，然后打开
\key{Open typewriter}。在电脑上选中文档输入位置，再点平板上的\key{Start}。
打字机始终以暂停状态打开。\key{Escape} 或\key{Stop} 暂停并释放按键。
离开打字机、切换终端、窗口失焦、打开原生应用或休眠时也会暂停。
断线重连后不会自动重发。

\begin{choices}
US / ASCII & 美式布局，关闭 Caps Lock。拒绝非 ASCII 文本。 \\
Mac & 在键盘输入源中选择 Unicode Hex Input，关闭 Caps Lock。
只支持到 U+FFFF，拒绝补充平面字符。 \\
Linux & 美式布局，关闭 Caps Lock。接收应用／输入法必须支持 Ctrl+Shift+U、
十六进制数字、Enter；并非每个应用都支持。 \\
Windows & 在 PC 上安装 WinCompose，将 Compose 设为 Right Alt，启用 Unicode 输入。
使用美式布局并关闭 Caps Lock。发送 Compose、u、六位十六进制数字、Enter。 \\
\end{choices}
从 \url{https://wincompose.info/} 获取 WinCompose。自定义规则不可覆盖此 Unicode 序列。
HID 没有通用的 Unicode 文本通道；结果取决于接收端配置和字体。先用短文本试验。

\key{Input method} 与终端共用输入引擎：Off／直接美式、Romaji、Pinyin、Zhuyin、
Wubi 86 和 US-International。只有已确认字符才会发送。更换输入法后需点 Start 继续。
选择\key{Off} 返回普通美式输入。\key{Unicode} 或 Alt+Space 打开的字符面板也向 USB 接收端输入。

Typewriter 是本地编辑器，Inkline 运行期间文档只保存在 RAM。普通输入会编辑本地文档；
开启实时发送后，每个已确认字符也会发送到连接的电脑。复制、剪切、粘贴、选择和双指滚动
都作用于本地编辑器。\key{Files \& send} 可用任意文件名加载或保存 UTF-8，
以每秒 5--80 个字符发送全文，并可停止发送。编辑时不会写入存储；只有 Save 或 Inkline
正常退出时才会写入。无标题文档在退出时使用隐藏恢复名
\path{~/.inkline-typewriter-draft}。\key{Exit typewriter} 返回 USB 设置。

\topic{发送文件与恢复网络}
安装器把 \code{inkline-usb} 和 \code{keyboard-send} 放入
\path{~/.local/bin}。\code{~/inkline usb}、\code{~/inkline type} 和
\code{inkline-type} 保留为兼容别名。Outpost 的独立工具保持原样。
\UsbExamples
UTF-8 文件\key{没有固定大小上限}。发送器分块读取，不在 RAM 中保存整个文档，
也不创建中转文件。可回读的普通文件会先完整验证再重新读取发送；发送时不要修改源文件。
管道按到达顺序验证，因此后续错误出现时前面的内容可能已经发出。
可选的 UTF-8 BOM 会被去除，CRLF／CR 转为换行。Enter 和 Tab 会操作当前应用。

\code{--cps}（默认 20）或 \code{--delay-ms} 设置速度；Unicode 需要更多按键报告。
\code{--start-delay} 留出切换输入焦点的时间。Dry-run 只打印报告，不发送 USB 按键。
Ctrl+C 或 \code{--stop} 取消并释放按键。一次只能有一个发送器。
交互队列满时会暂停；中断后重试前先检查接收文档。文本和临时状态不写入闪存。
接收端配置与其他设置一起在正常退出时保存。

点\key{Network} 或运行 \code{inkline-usb network} 恢复以太网。
Inkline 退出、崩溃或紧急重启时也会恢复网络。USB 模式不会保存为开机偏好。
\UsbRecoveryExample
定时器独立于 shell 运行。从网络模式切换出去应使用平板或 Wi-Fi SSH，不能使用 USB SSH。
切换失败会尝试恢复以太网。Inkline 不会抢占活动的卫星 PPP 或未知 USB 设备功能；
退出清理也遵循此规则。Opt+RightAlt+Backspace 紧急重启仍然可用，
但会关闭所有终端并丢失其中未保存的工作。

\topic{用快捷键启动程序}
在 Inkline shell 中注册快捷键。参数按字面传递；如需 shell 语法，请显式调用 shell。
快捷键程序通过 Bash 读取 \code{~/.bashrc}。

Settings → Command shortcuts（第 2 页）可为字母键指定命令。选择 Terminal 或 Native e-paper app，再按 Save。Remove 需确认后才删除；Reset draft 放弃编辑。T 已锁定。只有 Save 和确认删除会写入存储。命令行工具管理同一组快捷键。

所有全局快捷键现在使用 Opt+RightAlt，包括 T 和长按 Backspace 的紧急恢复组合。Folio 使用独立的 Opt 键；USB 键盘使用 Windows/Command（Super）键。普通 Ctrl+Alt 仍可用于 Emacs。升级后请退出并重新打开 Inkline；安装会保留已打开的终端。

\begin{shell}
~/inkline shortcut register e emacs
~/inkline shortcut register p python3
~/inkline shortcut list
~/inkline shortcut deregister p
\end{shell}

\key{Opt+RightAlt+E} 在空闲槽位打开 Emacs。
\code{~/inkline run PROGRAM ARG...} 可不注册直接启动。
重复注册会替换绑定；注销不影响正在运行的程序。
按键采用美式键盘的物理位置。标点请加引号，自定义程序请使用完整路径，并先安装程序。

对于自行绘制屏幕的原生 Qt／电子纸应用：
\begin{shell}
~/inkline shortcut register a --epaper /home/root/my-app
\end{shell}
一次运行一个原生应用，期间暂停 Inkline 及其 shell。启动器设置 Qt 电子纸环境与 RAM 缓存，
丢弃日志，并在应用退出后恢复 Inkline。应用应使用 Qt 输入，不要独占键盘设备。

\begin{notice}
\key{Opt+RightAlt+T 是固定快捷键，不能注册或注销。}
它会停止原生快捷键应用，并返回现有终端。
此组合不需要 Ctrl；Caps Lock 映射只在 Inkline 内部生效。
\end{notice}

\key{紧急恢复：}按住 Opt+RightAlt+Backspace 两秒，停止快捷键应用及其子进程并重新启动 Inkline。
所有终端会关闭，未保存的工作会丢失。以 root 运行的应用可能通过独占输入、修改服务或耗尽资源
使恢复失败；最后的办法是重启设备。

\topic{安装、升级与保留手册}
此安装包面向\key{运行固件3.27的 reMarkable 2}，使用系统自带 Qt 6.8 电子纸库。
最新下载与校验步骤见：

\InstallLink

安装器检查硬件、固件、字体、Qt 启动、RAM 临时目录和 swap 是否禁用。
版本保存在 \path{/home/root/.local/share/inkline} 下，键盘启动服务随系统启动。
更新不会关闭已有终端。\key{退出并重新打开 Inkline 才会使用新版本。}
旧版本保留用于回退，会继续占用存储空间。

安装包附带这一个包含九种语言的 PDF。笔记界面正在运行且 USB 网页接口已启用时，
安装器会通过该接口尝试导入\key{My files}。成功导入后记录校验值，避免重复导入同一手册。

如果自动导入不可用，请连接 USB、退出 Inkline，并在笔记界面的存储设置中启用
\key{USB web interface}，然后通过 SSH 运行：
\begin{shell}
~/inkline manual
\end{shell}
也可下载 PDF，通过 USB 浏览器界面或 reMarkable 桌面／移动应用导入。
导入器不会直接修改笔记数据库，也不会中断终端。
卸载 Inkline 后，已导入手册仍留在 My files。

通过 SSH 返回笔记或卸载 Inkline：
\begin{shell}
~/inkline stop
~/inkline uninstall
\end{shell}
可选软件（包括推荐的 LaTeX 环境）的安装说明保留在网站上。

\topic{工作原理、存储与 RAM}
Qt 将键盘事件交给共享映射器，处理快捷键与 Caps Lock。
libghostty-vt 编码终端按键，每个槽位通过独立伪终端（PTY）连接程序。
shell 获得 \code{HOME=/home/root}，交互式 Bash 读取 \code{~/.bashrc}。

程序输出经过终端解析器与 sixel 解码器。libghostty 管理字符单元、光标状态、备用屏幕、
回滚历史和 kitty 图像。渲染器保留灰度图，只重绘变化的行；系统 Qt 电子纸后端负责刷新面板。
独立的 evdev 守护进程处理全局快捷键和电源键，systemd 负责休眠与唤醒。

每个终端的核心分配上限为64 MiB；kitty 图像每个屏幕16 MiB；保留的 sixel 图像16 MiB。
渲染器和应用内存另计。空槽位不分配终端。九个大型应用可能超过可用 RAM。

临时图像、缓存和会话状态使用 RAM。已安装软件、保存的文档、退出时保存的设置以及应用写入的
文件使用持久存储。Inkline 不会禁止任意程序写文件。

可选 Python 包是标准 Python。为减少缓存写入，请在设备的 \code{~/.bashrc} 中添加：
\begin{shell}
export PYTHONDONTWRITEBYTECODE=1
export PIP_NO_CACHE_DIR=1
\end{shell}
运行 \code{source ~/.bashrc} 或打开新 shell。
使用 \code{python3 -m pip install --no-compile PACKAGE} 可避免 pip 在安装时单独进行的字节码编译。
Python 隔离模式会忽略 Python 环境变量。详细内容见网站的 Python 指南。

Inkline 采用 GPL-3.0-or-later 许可。安装包包含匹配的源代码归档与依赖许可说明。
测试记录区分自动检查与仍需实机确认的行为。
